\section{Contexte et Objectifs}

L'objectif de ce projet est de concevoir un robot qui depuis un point de départ, doit passer par des points de passage obligatoire et revenir à son point de départ le plus rapidement possible. Il devra tout au long du parcours suivre une ligne noire.

Le robot sera positionné au départ sur une pastille de couleur rouge, et devra donc passer par toutes les pastilles de couleur verte tout en empruntant le chemin le plus court. A chaque passage de pastille, le robot émettra un bip et devra afficher la couleur de la pastille sur le lcd. 

Pour réaliser à bien ce projet, nous avons mobilisé plusieurs compétences en électronique et en algorithmique, tout en appliquant les diverses connaissances vue en cours. De plus, nous avons du apprendre à travailler en équipe et à s'écouter les uns les autres.

Dans un premier temps, nous aborderons la partie électronique de notre robot en détaillant la conception et la réalisation. Nous expliquerons notre schéma électrique, nos choix de câblage ainsi que tous les choix techniques que nous avons pris. Nous réaliserons un bilan de puissance et une estimation de l'autonomie de notre robot. Puis nous inclurons une explication détaillée du rôle et du principe de fonctionnement des composants choisis. Un organigramme sera réalisé afin de détailler notre démarche. Pour finir, dans cette partie, nous expliquerons nos difficultés rencontrées et nos pistes d'améliorations.

Dans un second temps, nous aborderons la partie informatique où nous présenterons la mise en place du cycle en V. Comme nous avons vu dans nos cours d'algorithme, nous ferons dans un premier temps : l'analyse puis la conception préliminaire, la conception détaillée et pour finir nous traduirons les codes en C.
